% !PRINTED 319
\weberpage{319}
%<<CONT
$z = x + yi$, darin seien $\varphi$ und $\psi$ reelle Functionen von $x$ und
$y$, und es werde nun eine geschlossene Curve $f$ gezeichnet. Wir haben dann
den Satz:

\gesperrt{XI. Wenn die Curve $f$ durch die Curve $\varphi$ in $2m$ Segmente
getheilt wird, in denen $\varphi$ abwechselnd positiv und negativ ist, so ist
die Zahl der innerhalb $f$ liegenden Wurzeln von $F(z)$ höchstens gleich $m$;
und dasselbe gilt auch, wenn $f$ durch $\psi$ in $2m$ solcher Segmente
getheilt wird.}

Der Satz ist eine unmittelbare Folge des Charakteristikensatzes II,
\S.\,96; denn danach ist
\[
k = \tfrac{1}{2}[E(f;\varphi,\psi) - A(f;\varphi,\psi)];
\]
und wenn nun die Curve $f$ von der $\varphi$-Curve in $2m$ Punkten
geschnitten wird, so ist
\[
m = \tfrac{1}{2}[E(f;\varphi,\psi) + A(f;\varphi,\psi)],
\]
also
\[
k = m - A(f;\varphi,\psi).
\]

Darin ist aber nach dem erwähnten Satze $k$ gleich der Anzahl der im Inneren
von $f$ liegenden Wurzeln von $F$, und $A(f;\varphi,\psi)$ ist eine
jedenfalls nicht negative Zahl. Ganz ebenso kann man mit der Formel
\[
k = \tfrac{1}{2}[E(f;\psi,\varphi) - A(f;\psi,\varphi)]
\]
verfahren.

\begin{center}
\S.\,106.

Das Theorem von Rolle.
\end{center}

Wenn die Gleichung $f(x) = 0$ zwei reelle Wurzeln $\alpha$ und $\beta$ hat,
von denen die erste eine $a$-fache, die zweite eine $b$-fache ist, so können
wir
\weq{1}{f(x) = (x - \alpha)^a (x - \beta)^b f_1(x)}
setzen.

Ist nun $\alpha < \beta$ und liegt zwischen $\alpha$ und $\beta$ keine
Wurzel von $f(x) = 0$, so wird $f(x)$ und $f_1(x)$ in diesem ganzen
Intervall ein unverändertes Vorzeichen haben. Bilden wir von (1) die
Ableitung, so ergiebt sich
%CONT>>
